近期身分保真個人化生成方法雖能忠實重現參考人臉，卻普遍偏向正面視角；面對轉頭、回眸、抬頭低頭等姿態指令時，現有方法常崩潰回正面。本論文研究「指定姿勢個人化文字轉圖像生成」：給定單一參考身分影像與姿態指令，生成符合姿態之同一身分影像，並以不偏正面之指標評估。

本論文提出 Rotate-ID，建構於少步驟修正流擴散轉換器（FLUX.2-klein）之單參考影像個人化系統。上下文姿態控制驅動頭部轉動；透過重新分配骨幹對參考影像之注意力達成身分保真，並萃取為約200參數可學習閘門（$\gamma$-gate），於不增姿態代價下精煉身分保真。

人臉辨識器於正面可靠讀出身分，卻隨頭部轉離正面而低估相似度；CLIP對齊亦無法判斷姿態是否達成。本論文之姿態校正效能指標同時解決兩者：以可解釋三維姿態檢測驗證轉動、以姿態適配讀出器評分，並僅於達成姿態時計入身分分數，建立涵蓋8種姿態、40組請求之評測基準。

此指標下，Rotate-ID 偏航角準確率(82.6\%)大幅超越最強基準PhotoMaker-V2(36.1\%)，於過肩轉動達成無基準方法能及之角度，且為姿態感知身分保真下最佳之免調校方法。

\vspace{1em}
\noindent 關鍵字：個人化文字轉圖像生成、身分保真、頭部姿態控制、擴散轉換器、注意力重新分配、修正流、姿態校正評估
